\section*{Critical--diffuse normal form for high FEI ratio}

All logarithms are base two except \(\ln\).  Let \(f:\{-1,1\}^n\to
\{-1,1\}\) be nonconstant, let \(S\) have the spectral law
\(\Pr[S=A]=\widehat f(A)^2\), and put
\[
 H=H(S),\qquad I=\mathbb E|S|,\qquad R=H/I,
 \qquad \alpha_i=\Pr[i\in S]=I_i(f).
\]
For a relevant coordinate \(i\), write
\[
 \Delta_iH=H(f)-\tfrac12\bigl(H(f_{i=1})+H(f_{i=-1})\bigr).
\]
Call \(f\) restriction-critical when
\(\Delta_iH\ge R I_i(f)\) for every relevant \(i\).

The exact A003 restriction split is
\[
 \Delta_iH=E_i+J_i,
 \quad E_i=H(\mathbf1_{\{i\in S\}}\mid S\setminus\{i\}),
 \quad J_i=\operatorname{JS}(r_+,r_-),
\]
where \(r_+,r_-\) are the two conditional restriction spectra.  Its two
pointwise bounds are
\[
 E_i\le h_2(\alpha_i),\qquad
 J_i\le {2\over\ln2}\min\{\alpha_i,1-\alpha_i\}.
\]
Consequently every restriction-critical \(f\) obeys, coordinate by
coordinate,
\[
 R\le G(\alpha_i):={h_2(\alpha_i)\over\alpha_i}
   +{2\over\ln2}{\min\{\alpha_i,1-\alpha_i\}\over\alpha_i}. \tag{1}
\]

Set
\[
 R_0=G(1/2)=2+{2\over\ln2}.
\]
On \((0,1/2]\),
\[
 G(\alpha)={h_2(\alpha)\over\alpha}+{2\over\ln2},
 \qquad {d\over d\alpha}{h_2(\alpha)\over\alpha}
 ={\log_2(1-\alpha)\over\alpha^2}<0. \tag{2}
\]
Thus this branch is continuous and strictly decreasing from infinity to
\(R_0\).  Let \(\kappa:(R_0,\infty)\to(0,1/2)\) be its inverse.  If
\(\alpha\ge1/2\), then
\[
 G(\alpha)\le {1\over\alpha}+{2\over\ln2}\le R_0. \tag{3}
\]
Equations (1)--(3) prove that a restriction-critical \(f\) with \(R>R_0\)
has, for every relevant \(i\),
\[
 I_i(f)=\alpha_i\le\kappa(R)<1/2. \tag{4}
\]
Moreover, for \(0<\alpha\le1/2\),
\[
 {h_2(\alpha)\over\alpha}
 =\log_2{1\over\alpha}
  +{1-\alpha\over\alpha}\log_2{1\over1-\alpha}
 \le\log_2{1\over\alpha}+{1\over\ln2},
\]
because \(-\ln(1-\alpha)\le\alpha/(1-\alpha)\).  Hence
\[
 G(\alpha)\le\log_2{1\over\alpha}+{3\over\ln2}.
\]
Together with (1), this yields the explicit inversion
\[
 \boxed{\ I_i(f)\le e^3,2^{-R}\quad\hbox{for every relevant }i\ }.
 \tag{5}
\]

Combining (4)--(5) with the A001 critical-prime descent gives the following
normal form.  For every nonconstant genuine Boolean \(f\) with
\(R(f)>R_0\), there is a restriction descendant \(h\), tensor-prime after
dummy variables and constant factors are deleted, such that
\[
 R(h)\ge R(f),\qquad
 \max_i I_i(h)\le\kappa(R(h))\le e^3 2^{-R(h)}\le e^3 2^{-R(f)}. \tag{6}
\]
Therefore failure of FEI would force a genuine Boolean sequence \(h_k\)
that is simultaneously restriction-critical, tensor-prime,
\(R(h_k)\to\infty\), and exponentially diffuse in every coordinate.

Two immediate restricted consequences are:

\begin{enumerate}
\item For every fixed \(0<\delta\le1/2\), every restriction-critical
Boolean \(h\) with \(\max_i I_i(h)\ge\delta\) satisfies
\[
 {H(h)\over I(h)}\le G(\delta).
\]
Indeed, if the witnessing influence is at most \(1/2\), use monotonicity in
(2); if it exceeds \(1/2\), use (3) and \(G(\delta)\ge G(1/2)=R_0\).

\item If \(h\) in (6) has \(m\) relevant variables, then every positive
influence is a multiple of \(2^{1-m}\), and therefore
\[
 m\ge R(h)+1-{3\over\ln2}.
\]
\end{enumerate}

These statements do not prove FEI: the remaining class permits every
coordinate influence to vanish with dimension.  They sharpen the first
unsupported universal line from an unrestricted good-coordinate estimate to
an estimate on an explicit critical, tensor-prime, exponentially diffuse
Boolean class.
